%% Transportation Research Part C: Emerging Technologies
%% Elsevier article template (elsarticle)
\documentclass[preprint,12pt]{elsarticle}

\usepackage{hyperref}
\usepackage{amsmath}
\usepackage{booktabs}
\usepackage{graphicx}
\usepackage{multirow}
\usepackage{xcolor}
\usepackage{url}
\usepackage[T1]{fontenc}
\usepackage{listings}
\usepackage{caption}
\usepackage{subcaption}

\journal{Artificial Intelligence for Transportation}

\begin{document}

\begin{frontmatter}

\title{Spatial Grounding Failure in LLM-Generated Traffic Impact Analysis Recommendations: A Structured Evaluation Using Public Environmental Review Documents}

\author[1]{[Author]}
\address[1]{[Institution]}
\ead{[email]}

%% Abstract ────────────────────────────────────────────────────────────────────
\begin{abstract}
No structured evaluation exists for large language models (LLMs) applied to
Traffic Impact Analysis (TIA) mitigation generation, a high-cost, high-volume
deliverable in US transportation planning.
We address this gap with three contributions.
First, we construct a 33-record benchmark dataset by programmatically harvesting
transportation study appendices from two public environmental review portals
(California CEQAnet and Massachusetts MEPA), each record pairing a structured
project context with held-out ground-truth mitigation measures.
Second, we define a conditioned generation task that mirrors real practitioner
workflow: given a project description and the traffic analysis results (failing
intersections with Level of Service grades), generate the mitigation
recommendations a licensed engineer would prescribe.
Third, we evaluate five LLMs under four conditions (zero-shot, few-shot,
retrieval-augmented generation (RAG), and leave-one-out LoRA fine-tuning)
using a two-tier metric that separates \emph{intervention type} accuracy from
\emph{spatial location} accuracy.

The dominant finding is a systematic spatial grounding failure: all models
predict intervention categories with moderate competence
($F1_\text{type} = 0.21$--$0.30$) but fail to assign them to the correct
intersections ($F1_\text{loc} = 0.07$--$0.14$).
Few-shot prompting closes part of this gap and is the most reliable strategy
($+$13--29\% over zero-shot for most models).
LLM-as-judge plausibility scores (4.1--4.7/5) indicate that outputs are
professionally structured, suggesting a role as \emph{drafting aids} subject
to engineer review, rather than as autonomous decision-making tools.
\end{abstract}

\begin{keyword}
Traffic Impact Analysis \sep Large Language Models \sep Spatial Grounding \sep
Benchmark Evaluation \sep Transportation Planning \sep Few-Shot Learning \sep
Environmental Review \sep CEQA \sep Dataset
\end{keyword}

\end{frontmatter}

%% ─────────────────────────────────────────────────────────────────────────────
\section{Introduction}
\label{sec:intro}
%% ─────────────────────────────────────────────────────────────────────────────

Traffic Impact Analyses are among the most resource-intensive deliverables in
the US transportation planning process.
Required under state and local regulations for developments expected to generate
significant new vehicle trips, a typical TIA documents existing traffic conditions,
projects future demand using standardized trip generation rates
\citep{ITE2021}, evaluates intersection Level of Service (LOS) under build
conditions, and specifies physical and operational mitigation measures for any
identified deficiencies \citep{NCHRP2019}.
Preparation of a single TIA commonly requires 80--200 consultant hours, costing
\$15{,}000--\$60{,}000 depending on scope \citep[see, e.g.,][]{FHWA2014}.
With tens of thousands of TIAs filed annually across the United States, the
cumulative cost to the development and public-review process is substantial.

Scaled language models exhibit emergent competence in structured reasoning
and generation across multiple domains \citep{brown2020gpt3, wei2022emergent,
openai2023gpt4, bubeck2023sparks}, including clinical knowledge
\citep{MedicalLLM2023} and scientific information extraction
\citep{ScienceLLM2023}.
Transportation planning documents, with their distinctive combination of
spatial reasoning, engineering vocabulary, and regulatory constraints,
remain understudied.
Crucially, no prior work has asked whether LLMs can perform the specific TIA
step that consumes the most expert time: given traffic analysis results, what
physical improvements should be prescribed, and at which intersections?

This paper answers that question with three contributions:

\begin{itemize}
\item \textbf{A benchmark dataset and conditioned generation task.}
33 TIA records are programmatically harvested from two public environmental
review portals (California CEQAnet and Massachusetts MEPA) and paired with
held-out ground-truth mitigation measures extracted from real consultant reports.

\item \textbf{A two-tier evaluation framework separating intervention type
from spatial location.}
The metric independently assesses categorical coverage ($F1_\text{type}$)
and location-specific precision ($F1_\text{loc}$), exposing failure modes
that aggregate metrics would mask.

\item \textbf{An empirical finding: systematic spatial grounding failure.}
Across all five LLMs and four prompting conditions, models correctly
identify intervention categories but consistently misassign them to
incorrect intersections, suggesting spatial grounding in structured engineering
tasks remains an unresolved challenge.
\end{itemize}

The central finding is a \textbf{systematic spatial grounding failure}:
LLMs correctly identify intervention categories ($F1_\text{type}$ 0.21--0.30)
but consistently misassign them to incorrect intersections
($F1_\text{loc}$ 0.07--0.14).
This is not a general model error but a structured failure specific to spatial
reasoning in professional engineering tasks, and it persists across model sizes,
prompting strategies, and domain adaptation methods.
LLMs can support TIA drafting as assistants to licensed engineers, but
location-specific outputs require expert review before adoption.

The remainder of this paper is structured as follows.
Section~\ref{sec:background} reviews related work on TIA processes,
transportation document processing, and LLMs in engineering domains.
Section~\ref{sec:data} describes the dataset construction methodology.
Section~\ref{sec:method} defines the benchmark tasks and evaluation metrics.
Section~\ref{sec:results} presents experimental results.
Section~\ref{sec:discussion} discusses practical implications and limitations.
Section~\ref{sec:conclusion} concludes with directions for future work.

%% ─────────────────────────────────────────────────────────────────────────────
\section{Background and Related Work}
\label{sec:background}
%% ─────────────────────────────────────────────────────────────────────────────

\subsection{The Traffic Impact Analysis Process}

A Traffic Impact Analysis is a technical study quantifying the effect of a
proposed land development on the surrounding roadway network.
The standard methodology, codified by the Institute of Transportation Engineers
\citep{ITE2021} and adopted in varying forms by state departments of
transportation and local jurisdictions across the United States, proceeds in
five stages: (1) defining a study area and selecting study intersections based on
expected traffic influence; (2) collecting existing traffic counts and
characterizing baseline Level of Service (LOS) using the Highway Capacity Manual
\citep{HCM2022}; (3) generating project trip estimates using ITE
land-use-specific rates; (4) distributing and assigning project trips to the
network; and (5) re-evaluating LOS under projected build conditions and
identifying mitigation measures for intersections projected to fail the
jurisdiction's LOS standard (typically LOS~D or better).

Mitigation measures in a TIA fall into a limited taxonomy: physical capacity
improvements (turn lanes, road widening, signal installation), operational
improvements (signal timing optimization), access management (driveway
consolidation, turn restrictions), and demand management (TDM programs,
transit subsidies).
The specificity and defensibility of these recommendations, including naming the
exact intersection, lane geometry, and storage length, is a core deliverable
of the professional traffic engineer.

\subsection{Automation in Transportation Document Processing}

Prior work on NLP for transportation has primarily targeted operational data
sources: classifying traffic incidents, extracting entities from police crash
reports \citep{zhang2020secondary}, and analysing driver injury narratives.
Planning-level documents (TIAs, EIRs, corridor studies) have received
substantially less attention, because they are longer, more heterogeneous
in structure, and require domain-specific reasoning rather than named
entity extraction.

Recent work has begun exploring LLMs for transportation planning contexts.
\citet{yan2025llm} use LLM agents to simulate resident preferences in
transit policy referenda, and \citet{choi2026climate} apply RAG-based LLMs
to extract transportation policy items from municipal climate equity plans.
Neither study addresses the generation or evaluation of intersection-specific
mitigation measures from traffic analysis inputs.
This paper addresses that gap.

\subsection{Large Language Models for Technical Documents}

LLMs have shown strong performance on structured generation in several
professional domains, including clinical knowledge encoding
\citep{MedicalLLM2023} and structured information extraction from scientific
literature \citep{ScienceLLM2023}.
A recurring pattern is that outputs are \emph{categorically} appropriate
but lack the contextual specificity required for deployment without
human review.

Transportation planning sharpens this limitation.
A TIA recommendation is defensible only when it names the exact intersection,
specifies the physical change, and quantifies its dimensions.
Recommending signal timing optimization at the wrong intersection is not
merely unhelpful: it could contribute to unsafe or unworkable designs.
This specificity requirement motivates our separation of type accuracy
from location accuracy in evaluation.

\subsection{Retrieval-Augmented Generation}

Retrieval-Augmented Generation \citep{RAG2020, gao2023rag} extends LLM prompting
by retrieving semantically similar documents from a corpus and providing them
as in-context examples.
RAG has shown particular promise for specialized professional domains where
factual precision matters \citep{RAG2020, gao2023rag}, as examples from
similar cases can ground the model's outputs in domain conventions.

Few-shot prompting \citep{brown2020gpt3} provides complete input--output
demonstrations directly in the prompt.
\citet{min2022rethinking} show that the \emph{format} of demonstrations
matters more than their factual content, explaining why few-shot can be
effective even when retrieved examples come from different projects:
the model learns the expected output structure rather than the specific answers.
This distinction has implications for our benchmark: few-shot demonstrations
teach the model ``a recommendation looks like this JSON object with these
fields,'' while RAG provides relevant text fragments that may or may not
contain complete examples.

In the context of TIA generation, we hypothesize that few-shot will outperform
RAG because TIA mitigation language follows a highly formulaic structure that
can be learned from very few complete examples \citep{zhao2021calibrate}.

%% ─────────────────────────────────────────────────────────────────────────────
\section{Dataset Construction}
\label{sec:data}
%% ─────────────────────────────────────────────────────────────────────────────

\subsection{Source Selection}

We identified two publicly accessible repositories of transportation impact
study documents that expose programmatic interfaces:

\textbf{California CEQA State Clearinghouse (CEQAnet).}
California's Environmental Quality Act (CEQA) requires that all Draft and
Final Environmental Impact Reports (DEIRs/FEIRs) for projects exceeding
defined thresholds be submitted to the State Clearinghouse at
\url{ceqanet.opr.ca.gov}.
Since approximately 2015, the portal has accepted electronic document
submissions as separate PDF attachments indexed by project and submittal type.
Critically, many projects include standalone Transportation Impact Analysis
appendices (labeled as ``Appendix K -- Transportation'', ``Traffic Impact
Study'', etc.) rather than embedding the traffic analysis within the main EIR
narrative, enabling extraction of self-contained TIA documents.

\textbf{Massachusetts MEPA e-Monitor.}
Massachusetts's Environmental Policy Act (MEPA) requires environmental
notifications and impact reports for major development projects, including a
transportation impact assessment.
The MEPA e-Monitor at \url{eeaonline.eea.state.ma.us} exposes a REST API
(identified through browser network traffic inspection) that provides
programmatic access to project metadata, submission records, and attached
PDF documents.

Both portals provide complete document provenance (project name, lead agency,
municipality, submission date), enabling reliable metadata construction for
each record in the dataset.
Full dataset documentation following the Datasheets for Datasets framework
\citep{gebru2021datasheets} is provided in Appendix~\ref{app:datasheet}.

\subsection{Programmatic Document Acquisition}

\subsubsection{CEQAnet Acquisition Pipeline}

The CEQAnet portal renders project pages as server-side HTML, enabling direct
scraping without browser automation.
We query the daily filing index at
\texttt{ceqanet.opr.ca.gov/Search?ReceivedDate=\{YYYY-MM-DD\}} for each
month-end and mid-month date in the range 2015--2019, focusing on this
pre-2023 window because California Senate Bill 743 (effective 2023 for most
jurisdictions) shifted CEQA transportation analysis from Level of Service
to Vehicle Miles Traveled (VMT) metrics, rendering post-2022 transportation
chapters methodologically incompatible with our LOS-based evaluation schema.

For each date, we identify rows with document type ``EIR'' and retrieve the
project page.
We then extract all PDF attachment records and filter for documents whose
filename or label contains transportation-specific keywords
(``Traffic Impact'', ``Transportation Analysis'', ``Level of Service'',
``Trip Generation'').
We apply a size filter (1--30 MB) to exclude cover pages and raw data
appendices, and download the resulting candidate PDFs.

\subsubsection{MEPA API Acquisition Pipeline}

Through reverse engineering of the MEPA e-Monitor's Angular frontend, we
identified a public AWS API Gateway endpoint at
\texttt{t7i6mic1h4.execute-api.us-east-1.amazonaws.com}
that accepts requests authenticated via a static API key embedded in the
application bundle.
We enumerate all published Environmental Monitor editions from 2020--2023
using the \texttt{/api/Publishing/publication} endpoint,
retrieve project submission lists for each edition via
\texttt{/api/PublicationHistory}, and download transportation study attachments
via \texttt{/api/Attachment/\{fileServiceId\}}.
A session warm-up step (loading the MEPA portal) is required to establish an
IP-based server session that authorizes file downloads.

\subsection{PDF Parsing and Section Extraction}

Downloaded PDFs are processed using \texttt{pdfplumber} \citep{pdfplumber}.
We apply a file-size-based page cap (all pages for files $<$5 MB; up to 200
pages for files 5--15 MB; up to 150 pages for larger files) to manage
processing time on multi-chapter EIR volumes.

Section boundaries are detected using a two-tier pattern set: standard TIA
section headers (``Existing Conditions'', ``Trip Generation'',
``Mitigation Measures'', etc.) and EIR-specific equivalents
(``IV.D Transportation and Circulation'',
``Intersection Level of Service Analysis'', etc.).
Each detected section is assigned one of three roles:

\begin{itemize}
\item \textbf{Input}: Sections 1--8 of the TIA sequence (project description,
study area, existing conditions, trip generation, trip distribution,
traffic assignment, future conditions analysis). These constitute the
model's input at evaluation time.
\item \textbf{Ground truth}: The findings and mitigation sections of the
report. These are held out and used only for evaluation.
\item \textbf{Other}: Preamble, appendices, and administrative sections
excluded from both roles.
\end{itemize}

\subsection{Structured Information Extraction}

Input and ground-truth sections are converted to structured JSON using
Cloudflare Workers AI (\texttt{@cf/meta/llama-3.3-70b-instruct-fp8-fast})
via two targeted prompts:
an \emph{input extraction prompt} that maps project description text to a
structured schema (project type, location context, trip generation summary,
study intersection list with LOS grades), and a \emph{ground-truth extraction
prompt} that extracts issues (deficient intersections with LOS grade,
V/C ratio, and deficiency type) and recommendations (mitigation measures
with location, measure type, description, timing, and responsible party).

Records are retained in the final dataset only if:
(1) both input and ground-truth sections are detected;
(2) the input text contains at least 8{,}000 characters;
(3) the ground-truth text contains at least 5{,}000 characters; and
(4) the extracted ground truth contains at least one recommendation.
A leakage check verifies that no finding-specific language
(e.g., ``operates at LOS E'', ``dedicated left-turn lane'')
appears in the input fields.

\subsection{Dataset Characteristics}

The resulting dataset comprises \textbf{33 records} from \textbf{23
California} and \textbf{10 Massachusetts} projects spanning
\textbf{2012--2024}.
Table~\ref{tab:dataset} summarizes key statistics.

\begin{table}[h]
\centering
\caption{Dataset statistics.}
\label{tab:dataset}
\begin{tabular}{lr}
\toprule
Statistic & Value \\
\midrule
Total records & 33 \\
Records with both issues and recommendations & 17 \\
California (CEQAnet) & 23 \\
Massachusetts (MEPA) & 10 \\
Year range & 2012--2024 \\
Unique project types & 8 \\
Total ground-truth issues & 91 \\
Total ground-truth recommendations & 189 \\
Mean recommendations per case & 5.73 \\
Mean input text length (chars) & 68,198 \\
Mean GT text length (chars) & 67,837 \\
Mean chunks per case & 86.0 \\
\bottomrule
\end{tabular}
\end{table}

Project types include mixed-use development ($39\%$), industrial/warehouse
($21\%$), residential multifamily ($6\%$), commercial/office
($12\%$), and institutional uses ($9\%$).
The dataset covers 16 California cities and municipalities across
10 counties, plus 8 Massachusetts municipalities.

Each record in the final dataset is stored as a JSON object with:
(1) an \texttt{input} field containing project type, location context,
description, and known conditions (study intersections with baseline LOS);
(2) a \texttt{ground\_truth} field with structured lists of
\texttt{issues[]} and \texttt{recommendations[]};
(3) a \texttt{metadata} field with source URL, agency, year, and an
\texttt{exclude\_from\_retrieval} list enabling leave-one-out evaluation.

%% ─────────────────────────────────────────────────────────────────────────────
\section{Benchmark Design and Evaluation}
\label{sec:method}
%% ─────────────────────────────────────────────────────────────────────────────

\subsection{Task Definition}

We define a single primary task, \emph{conditioned recommendation generation},
evaluated under four experimental conditions.

\textbf{Task: Conditioned Mitigation Generation.}
Given (1) the structured project input (project type, location context,
description, trip generation summary, existing intersection LOS) and (2) the
traffic analysis results (a structured list of failing intersections with
LOS grade, V/C ratio, and deficiency type), generate the mitigation
recommendations a licensed traffic engineer would prescribe.

This task design mirrors actual TIA practice: the quantitative LOS analysis
is completed by software (Synchro, HCM); the engineer then writes the
mitigation section. We evaluate whether LLMs can perform the latter step
when given the former's output.
Evaluated on the $n = 17$ cases that have both non-empty issues and
recommendations in the ground truth.

\textbf{Condition 1 (Zero-shot):} Project input + failing intersections only.
The model must determine appropriate fixes with no examples.

\textbf{Condition 2 (Few-shot):} Same input plus 2--3 complete demonstrations
selected from the most similar other cases (by project type and state),
formatted as complete (input, analysis, recommendations) triplets.
Unlike RAG, few-shot provides fully-worked examples rather than retrieved fragments.

\textbf{Condition 3 (RAG):} Same input plus the top-$k = 8$ text chunks
retrieved by cosine similarity over BGE-M3 embeddings \citep{BGEM3},
with leave-one-out exclusion of the query case.

\textbf{Condition 4 (LoRA fine-tuning):} Leave-one-out fine-tuning of
Mistral-7B-Instruct-v0.2 \citep{jiang2023mistral} using rank-8 LoRA
\citep{hu2022lora} with 4-bit NF4 quantization \citep{dettmers2023qlora}.
Each of the 17 LOO folds trains on 16 examples and evaluates on the
held-out case (full hyperparameters in Appendix~\ref{app:compute}).
This condition tests whether weight-based adaptation outperforms
in-context learning at this dataset scale.

\subsection{Models Evaluated}

We evaluate five LLMs served via Cloudflare Workers AI, one fine-tuned
adapter, and a frequency baseline:

\begin{itemize}
\item \textbf{Llama-3.3-70B} \citep{touvron2023llama2} (\texttt{@cf/meta/llama-3.3-70b-instruct-fp8-fast}), 70B instruction-tuned
\item \textbf{GPT-OSS-120B} \citep{openai2023gpt4} (\texttt{@cf/openai/gpt-oss-120b}), 120B
\item \textbf{Nemotron-3-120B} (\texttt{@cf/nvidia/nemotron-3-120b-a12b}), 120B
\item \textbf{Qwen3-30B} (\texttt{@cf/qwen/qwen3-30b-a3b-fp8}), 30B extended-thinking model
\item \textbf{Gemma-3-12B} \citep{team2024gemma} (\texttt{@cf/google/gemma-3-12b-it}), 12B
\item \textbf{Mistral-7B LoRA} \citep{jiang2023mistral} (\texttt{@cf/mistral/mistral-7b-instruct-v0.2-lora}), fine-tuned via LOO LoRA \citep{hu2022lora} using 4-bit NF4 quantization \citep{dettmers2023qlora} (Condition 4 only)
\item \textbf{Frequency Baseline}: predicts the most common measure types
observed in LOO training records, calibrated by project type.
Does not use the provided failing intersections.
\end{itemize}

All LLMs for Conditions 1--3 share a system instruction establishing the
transportation engineering context and requesting JSON-only output;
temperature is set to 0.1. Qwen3-30B uses 4,096 output tokens to
accommodate its extended chain-of-thought reasoning prior to the JSON response.

\subsection{Evaluation Metrics}

\subsubsection{Recommendation Evaluation}

We evaluate recommendations using a \emph{two-tier matching} approach
that mirrors professional TIA review practice:

\textbf{Location-aware F1 (primary metric, $F1_\text{loc}$).}
A predicted recommendation is considered a true positive if:
(a) the measure type exactly matches a ground-truth recommendation; \emph{and}
(b) the predicted location is sufficiently similar to the GT location
(Jaccard similarity $\geq 0.5$ on meaningful word tokens after stopword removal).
When either the predicted or ground-truth location is unspecified (e.g.,
project-wide TDM programs), only measure-type agreement is required.
Precision, recall, and F1 are computed using greedy one-to-one matching.

\textbf{Type-only F1 (secondary metric, $F1_\text{type}$).}
The multi-set intersection of predicted and ground-truth measure types,
regardless of location. This provides a ceiling on $F1_\text{loc}$ and
measures categorical coverage.

The gap between $F1_\text{type}$ and $F1_\text{loc}$ quantifies location
specificity failure: models know \emph{what} to recommend but not
\emph{where}.

\subsubsection{LLM-as-Judge Evaluation}

Automatic metrics may not fully capture whether outputs are professionally
usable. We supplement them with an LLM-as-judge evaluation
\citep{LLMJudge2023}, a methodology increasingly used to assess open-ended
generation quality \citep{liang2022helm, chang2024survey}: for each
prediction, we prompt
Llama-3.3-70B (acting as a transportation engineering reviewer) to score
the output on three 1--5 scales:
\begin{itemize}
\item \textbf{Domain plausibility}: Would a licensed traffic engineer write
recommendations like these?
\item \textbf{Location specificity}: Are the named locations plausible given
the project context?
\item \textbf{GT alignment}: Do the predicted measure types match what the
ground truth recommends?
\end{itemize}
An overall score is also requested. This provides a practitioner-oriented
complement to the automatic metrics.

\subsection{Leave-One-Out Evaluation}

All four conditions use leave-one-out (LOO) evaluation: for any given case,
the retrieval corpus (Conditions 2--3), the frequency baseline training set,
and the fine-tuning training set (Condition 4) all exclude that case.
This ensures no data leakage while using all available data for both
training and evaluation.

%% ─────────────────────────────────────────────────────────────────────────────
\section{Results}
\label{sec:results}
%% ─────────────────────────────────────────────────────────────────────────────

Results are organised around the paper's central thesis: that LLMs show
categorical competence but spatial grounding failure.
We first characterise the ground-truth distribution (Section~\ref{sec:results}),
which is required to interpret the primary metric, then present the
four-condition comparison, the type-vs-location gap analysis,
and qualitative error examples that concretise the failure mode.

\subsection{Ground-Truth Type Distribution}

We first characterise the ground-truth recommendation distribution,
which directly affects evaluation interpretation.
Of the 71 ground-truth recommendations across the 17 Task C cases,
36 (51\%) are intersection-specific (a named location with two or more
meaningful road-name tokens) and 35 (49\%) are project-wide
(``Future developments'', ``Project site'', ``not specified'', or similar).
Intersection-specific recommendations cluster around physical capacity
improvements (\texttt{add\_turn\_lane}: 8/9 specific;
\texttt{new\_traffic\_signal}: 5/5 specific;
\texttt{signal\_timing\_optimization}: 4/5 specific),
while project-wide recommendations dominate TDM programs
(\texttt{tdm\_measure}: 5/5 project-wide),
access requirements (\texttt{access\_modification}: 8/12 project-wide),
and road obligations (\texttt{road\_widening}: 7/9 project-wide).

This distribution explains the \emph{specificity paradox} discussed in
Section~\ref{sec:discussion}: the frequency baseline, which always
predicts unspecified locations, matches project-wide GT items via
the Tier~2 (type-only) fallback and scores $F1_\text{loc} = 0.249$,
numerically higher than all LLMs.
This is a metric artefact, not a genuine performance advantage, as
confirmed by LLM-as-judge scores (Frequency Baseline: 2.00/5 vs.\ LLMs: 3.35--3.47/5).

\subsection{Four-Condition Comparison}

Table~\ref{tab:main_results} presents $F1_\text{loc}$ results for all
models across all four conditions ($n = 17$ cases, conditioned task).

\begin{table}[h]
\centering
\caption{Conditioned recommendation generation: $F1_\text{loc}$ by model and condition ($n = 17$).
Bolded values indicate the best condition per model.
Frequency baseline cannot perform few-shot or fine-tuning; RAG result is identical to zero-shot.}
\label{tab:main_results}
\begin{tabular}{lrrrr}
\toprule
Model & Zero-shot & Few-shot & RAG & Fine-tuned (LOO) \\
\midrule
Qwen3-30B          & 0.109 & \textbf{0.123} & 0.086 & --- \\
Llama-3.3-70B      & 0.096 & \textbf{0.124} & 0.070 & --- \\
Gemma-3-12B        & 0.097 & 0.118 & \textbf{0.141} & --- \\
GPT-OSS-120B       & \textbf{0.090} & 0.084 & 0.048 & --- \\
Nemotron-120B      & 0.070 & 0.053 & \textbf{0.072} & --- \\
Mistral-7B LoRA    & --- & --- & --- & 0.050 \\
Frequency Baseline & 0.249$^\dagger$ & --- & 0.249$^\dagger$ & --- \\
\midrule
\multicolumn{5}{l}{$^\dagger$Metric artefact: matches project-wide GT via Tier~2 fallback; judge score 2.00/5.} \\
\bottomrule
\end{tabular}
\end{table}

\textbf{Few-shot is the most consistent improvement.}
Few-shot prompting improves $F1_\text{loc}$ for 3 of 5 models
(Llama: $+$29\%, Qwen3: $+$13\%, Gemma: $+$22\%)
and ties for one (Gemma's best is RAG but few-shot is close).
Only Nemotron and GPT-OSS decline under few-shot.
The improvement is consistent across project types and geographic regions.

\textbf{RAG is heterogeneous.}
RAG provides Gemma's best result (0.141, $+$45\% over zero-shot)
and modestly helps Nemotron, but degrades performance for Llama ($-$27\%),
GPT-OSS ($-$47\%), and Qwen3 ($-$21\%).
We attribute this to model-dependent retrieval behaviour: Gemma appears
to use retrieved examples as structural templates, while larger models
may over-anchor to retrieved intersection names from unrelated projects.

\textbf{LoRA fine-tuning underperforms zero-shot.}
The LOO-fine-tuned Mistral-7B achieves $F1_\text{loc} = 0.050$, below
every zero-shot LLM (range 0.070--0.109).
Table~\ref{tab:finetuned_outputs} summarises the distribution of
fine-tuned outputs: 7/17 folds produced empty predictions;
6/17 predicted \texttt{no\_mitigation\_required} regardless of the
failing intersections provided; only 4/17 predicted \texttt{tdm\_measure}.
The model collapsed to two safe defaults, never predicting
\texttt{add\_turn\_lane}, \texttt{new\_traffic\_signal}, or
\texttt{signal\_timing\_optimization}, which together constitute
19 of 71 GT recommendations (27\%).
This failure mode is consistent with catastrophic overfitting on 16
training examples: the model memorised surface response patterns at
the expense of its zero-shot generalization ability.

\begin{table}[h]
\centering
\caption{LOO fine-tuned Mistral-7B prediction distribution vs.\ GT ($n=17$ folds).}
\label{tab:finetuned_outputs}
\begin{tabular}{lrr}
\toprule
Output type & Fine-tuned count & GT count \\
\midrule
Empty prediction                & 7  & 0 \\
\texttt{no\_mitigation\_required} & 6  & 1 \\
\texttt{tdm\_measure}             & 4  & 5 \\
\texttt{add\_turn\_lane}          & 0  & 9 \\
\texttt{signal\_timing\_optimization} & 0 & 5 \\
\texttt{new\_traffic\_signal}     & 0  & 5 \\
Other types                     & 0  & 46 \\
\bottomrule
\end{tabular}
\end{table}

\subsection{Type-Only vs.\ Location-Aware F1}

Table~\ref{tab:type_vs_loc} presents both metrics for the zero-shot condition,
illustrating the categorical/spatial gap across models.

\begin{table}[h]
\centering
\caption{Zero-shot condition: location-aware vs.\ type-only F1 ($n=17$).
The gap quantifies spatial grounding failure.}
\label{tab:type_vs_loc}
\begin{tabular}{lrrr}
\toprule
Model & $F1_\text{loc}$ & $F1_\text{type}$ & Gap ($F1_\text{type} - F1_\text{loc}$) \\
\midrule
Llama-3.3-70B  & 0.096 & 0.298 & 0.202 \\
Qwen3-30B      & 0.109 & 0.259 & 0.150 \\
GPT-OSS-120B   & 0.090 & 0.242 & 0.152 \\
Nemotron-120B  & 0.070 & 0.209 & 0.139 \\
Gemma-3-12B    & 0.097 & 0.222 & 0.125 \\
\bottomrule
\end{tabular}
\end{table}

The gap is substantial and consistent across all models (0.13--0.20 F1 points).
Models correctly anticipate what type of improvement is needed but cannot
reliably assign it to the correct intersection.
This is the central finding of the paper: \textbf{LLMs exhibit semantic
competence but spatial grounding failure.}

\subsection{Qualitative Error Analysis}

Table~\ref{tab:qualitative} illustrates the three recurring failure modes
across three representative cases.

\begin{table*}[ht]
\centering
\caption{Representative prediction failures (zero-shot, Llama-3.3-70B).
Each row shows a failing intersection given as input, the ground-truth
mitigation, and the model's prediction.}
\label{tab:qualitative}
\begin{tabular}{p{3.8cm}p{4.2cm}p{4.2cm}p{2.4cm}}
\toprule
Case / Input intersection & Ground truth & Model prediction & Failure type \\
\midrule
\textbf{CEQ-CA-010282} (mixed-use, Santa Clara)
Lafayette St \& El Camino Real (LOS D) &
``Driveways of future developments should meet minimum width requirements''
(\texttt{access\_modification}, project-wide) &
``Optimize signal timing at Lafayette St \& El Camino Real''
(\texttt{signal\_timing\_optimization}, same location) &
Wrong type at correct location \\
\addlinespace
\textbf{CEQ-CA-080060} (industrial, Perris)
Harvill Av. \& Perry St. (LOS F) &
``Install a stop control and 100-ft WB left turn lane at Driveway 1''
(\texttt{access\_modification}, Driveway 1) &
``Add dedicated EBL lane on Harvill Av. at Perry St.''
(\texttt{add\_turn\_lane}, correct main intersection) &
Correct type, wrong sub-location \\
\addlinespace
\textbf{MEPA-MA-16468} (TOD, Cambridge)
22 failing intersections incl.\ Binney St./Third St. (LOS F) &
``Construct sidewalk-level bicycle lanes on Potter Street''
(\texttt{pedestrian\_improvement}, non-failing street) &
``Install traffic signal at Binney St./Third St.''
(\texttt{new\_traffic\_signal}, failing intersection) &
Correct location, project type mismatch \\
\bottomrule
\end{tabular}
\end{table*}

Three patterns emerge.
\textbf{Wrong type at correct location} (Row 1): the model correctly names
the failing intersection but prescribes a standard engineering fix
(signal retiming) when the GT calls for a planning-level design obligation
(driveway standards). The model lacks context about the EIR's planning scope.
\textbf{Correct type, wrong sub-location} (Row 2): the model identifies
the right class of improvement (turn lane) at the right arterial but
targets the signalized intersection rather than the site driveway. Sub-location
granularity within a corridor is systematically lost.
\textbf{Project type mismatch} (Row 3): for a large transit-oriented
development, the GT prescribes multimodal corridor improvements;
the model defaults to intersection engineering. This reflects a bias
toward the most common training patterns (signal and turn-lane fixes)
regardless of project type.

\subsection{LLM-as-Judge Evaluation}

Table~\ref{tab:judge} presents LLM-as-judge scores for the zero-shot condition,
which reveal that professional plausibility is substantially higher than
automatic metrics suggest.

\begin{table}[h]
\centering
\caption{LLM-as-judge evaluation (zero-shot condition, 1--5 scale, $n=17$).
DP: domain plausibility. LS: location specificity. GT-A: GT alignment.}
\label{tab:judge}
\begin{tabular}{lrrrr}
\toprule
Model & DP & LS & GT-A & Overall \\
\midrule
Nemotron-120B      & 4.70 & 4.24 & 2.85 & 3.64 \\
Gemma-3-12B        & 4.47 & 3.76 & 2.70 & 3.47 \\
Qwen3-30B          & 4.18 & 3.76 & 2.21 & 3.47 \\
GPT-OSS-120B       & 4.53 & 3.59 & 2.45 & 3.41 \\
Llama-3.3-70B      & 4.06 & 3.59 & 2.27 & 3.41 \\
Frequency Baseline & 2.85 & 1.00 & 1.76 & 1.88 \\
\bottomrule
\end{tabular}
\end{table}

Domain plausibility scores (4.06--4.70) are high across all LLMs,
indicating that outputs use correct professional vocabulary and formatting.
Location specificity (3.59--4.24) is consistently the weakest dimension,
corroborating the $F1_\text{loc}$ results.
Crucially, the judge ranks LLMs (3.41--3.64) far above the frequency
baseline (1.88), confirming that the baseline's numerical advantage
in $F1_\text{loc}$ does not reflect practitioner utility.

%% ─────────────────────────────────────────────────────────────────────────────
\section{Discussion}
\label{sec:discussion}
%% ─────────────────────────────────────────────────────────────────────────────

\subsection{Primary Findings}

Four findings emerge consistently across conditions and evaluation methods.

\textbf{Finding 1: Semantic competence, spatial grounding failure.}
The gap between $F1_\text{type}$ ($0.21$--$0.30$) and $F1_\text{loc}$
($0.07$--$0.14$) is the central result of this study.
Models correctly identify what category of mitigation is appropriate for a
given project type: they know that a failing signalized intersection near
a warehouse development likely needs a turn lane, and that a mixed-use
development in a transit corridor likely needs a TDM program.
What they cannot do is assign those interventions to the correct specific
intersection.
This is a \emph{spatial} failure, not a \emph{semantic} one.
LLMs retain strong categorical knowledge but cannot reliably anchor it to
specific physical locations, consistent with documented limitations in
geographic and spatial reasoning \citep{bubeck2023sparks, openai2023gpt4}.
Importantly, this failure is \emph{structured}: it manifests as three
reproducible error modes (Table~\ref{tab:qualitative}) rather than random
noise, suggesting spatial grounding in professional engineering tasks is a
distinct, unresolved challenge for current LLMs.
Outputs are mislocated, not nonsensical.

\textbf{Finding 2: Few-shot prompting is the most sample-efficient adaptation strategy.}
Providing 2--3 complete demonstrations of the task consistently improves
performance for most models (3 of 5, gains of $+$13\%--$+$29\%),
outperforming both RAG and LoRA fine-tuning.
The explanation is structural: few-shot demonstrations show the model
exactly what a completed mitigation section looks like for a comparable
project, teaching format and vocabulary without the instability of
weight updates on tiny datasets.

\textbf{Finding 3: LoRA fine-tuning with 16 examples degrades performance.}
The LOO-fine-tuned Mistral-7B \citep{jiang2023mistral} using QLoRA
\citep{dettmers2023qlora} achieves $F1_\text{loc} = 0.050$, below
every zero-shot LLM and below its own estimated zero-shot baseline.
The model collapsed to two outputs: empty predictions (7/17 folds) and
\texttt{no\_mitigation\_required} (6/17 folds).
This behaviour reflects a documented failure mode of fine-tuning on very
small datasets \citep{zhao2021calibrate}: the model memorises surface-level
patterns from the instruction/response format while losing the generalisation
ability established during pretraining, sometimes called ``catastrophic
forgetting'' in the continual learning literature.
The finding that in-context few-shot learning \citep{brown2020gpt3} with
2--3 examples outperforms weight-based adaptation with 16 examples adds
to a growing body of evidence that the crossover point between prompting
and fine-tuning depends strongly on dataset scale and domain specificity
\citep{min2022rethinking}.
For this task, the threshold is somewhere above 16 examples; future work
with a larger dataset should establish it more precisely.

\textbf{Finding 4: The specificity paradox and its resolution.}
The frequency baseline achieves $F1_\text{loc} = 0.249$, numerically
higher than all LLMs, yet the judge rates it 1.88/5 vs.\ 3.41--3.64
for LLMs.
This paradox arises because 49\% of GT recommendations are project-wide
(no intersection named), which our Tier~2 fallback matches on type alone.
The baseline, which never names locations, exploits this fallback
disproportionately.
LLMs that attempt intersection-specific recommendations are evaluated
under the stricter Tier~1 criterion, penalising the more useful behaviour.
This finding has implications for evaluation design in spatial AI
systems more broadly: metrics requiring exact location agreement
may penalise specificity over abstention.

\subsection{Practical Implications for Transportation Practice}

The results support a specific, bounded role for LLMs in TIA practice:
\textbf{drafting support, not decision support.}
LLMs can assist engineers in generating initial mitigation language by predicting the category of improvement needed (signal timing, turn lane,
TDM program) with reasonable accuracy.
However, the spatial grounding failure documented here means that
location-specific recommendations require mandatory engineer review before
adoption.
Current outputs should not be used for unreviewed insertion into final TIA reports or autonomous decision-making.

The most defensible near-term application is \emph{consistency checking}:
flagging TIA reports that omit commonly required mitigation categories for
comparable project types, for review by agency staff.
This exploits LLM categorical competence while bypassing the location
specificity required for engineering precision.

\subsection{Limitations}

This study is a \emph{pilot benchmark} intended to establish methodology and
baselines. Four limitations should be noted explicitly.

\textbf{Small dataset (33 records, 17 Task C cases).}
Statistical variance is high at this scale.
Per-model F1 differences of $<0.02$ should not be interpreted as meaningful.
The primary value of the current results is directional, not definitive.
Expanding to 200+ records is tractable with the provided pipeline and is
the most important next step.

\textbf{LLM-extracted ground truth.}
Recommendations were extracted from source PDFs using Cloudflare Workers AI,
not validated by licensed traffic engineers.
Extraction errors inflate false-negative rates; the true model performance
may be somewhat higher than reported.
Future work should include expert spot-checking of at least 15--20 records.

\textbf{Geographic scope (California and Massachusetts only).}
Both states have unusually comprehensive public environmental review
requirements; their TIA structure and vocabulary may differ from states
without CEQA or MEPA equivalents.
Generalizability to the broader US practitioner landscape requires replication
across additional states.

\textbf{Post-2022 methodology shift.}
California's SB 743 shifted CEQA analysis from LOS to Vehicle Miles Traveled
(VMT). This benchmark covers only LOS-based practice (pre-2023 California,
all Massachusetts records), which remains current in most US jurisdictions.
VMT-based mitigation generation is a distinct and open research problem.

%% ─────────────────────────────────────────────────────────────────────────────
\section{Conclusion}
\label{sec:conclusion}
%% ─────────────────────────────────────────────────────────────────────────────

This paper presents the first structured evaluation of LLMs on conditioned
TIA mitigation generation.
The core finding is clear: \textbf{LLMs exhibit semantic competence but
spatial grounding failure.}
They predict intervention categories at a reasonable rate
($F1_\text{type}$ 0.21--0.30) but consistently misassign those interventions
to incorrect intersections ($F1_\text{loc}$ 0.07--0.14).
This gap persists across model sizes, architectures, and prompting strategies.
Qualitative analysis identifies three failure modes: wrong type at correct
location, correct type at wrong sub-location, and project-type mismatch
defaulting to generic engineering fixes.

Few-shot prompting with 2--3 complete demonstrations is the most reliable
strategy, providing consistent gains for most models.
LoRA fine-tuning with 16 examples degrades performance, establishing that
the current dataset is below the threshold needed for effective weight-based
adaptation.
LLMs can support TIA drafting as assistants to licensed engineers, but the
spatial grounding failure documented here means location-specific outputs
require mandatory expert review before use.

The most important direction for future work is dataset expansion.
The pipeline released alongside this paper can harvest hundreds more records
from CEQAnet and MEPA, and extending to VMT-based analysis is an open problem.
Scaling the dataset will establish whether spatial grounding failure is a
fundamental limitation of current LLMs on this task, or a phenomenon that
fine-tuning with sufficient domain examples can overcome.

The dataset construction pipeline and all benchmark scripts are available
at \url{https://github.com/[author]/tia-llm-benchmark}.

%% ─────────────────────────────────────────────────────────────────────────────
\section*{CRediT Author Statement}
%% ─────────────────────────────────────────────────────────────────────────────

\textbf{[Author]}: Conceptualization, Methodology, Software, Data Curation,
Formal Analysis, Writing -- Original Draft, Writing -- Review \& Editing.

\section*{Declaration of Competing Interests}

The author declares no competing interests.

\section*{Data Availability}

The dataset, pipeline code, and benchmark scripts are publicly available at
\url{https://github.com/[author]/tia-llm-benchmark}.
Source documents are publicly available from the California Governor's Office
of Planning and Research (\url{ceqanet.opr.ca.gov}) and the Massachusetts
Executive Office of Energy and Environmental Affairs
(\url{eeaonline.eea.state.ma.us}).

%% ─────────────────────────────────────────────────────────────────────────────
\appendix
%% ─────────────────────────────────────────────────────────────────────────────

\section{Compute Resources and Reproducibility}
\label{app:compute}

\subsection{Hardware Environment}

All inference experiments (Conditions 1--3) were run via the
Cloudflare Workers AI API, which served models on shared GPU infrastructure.
No local GPU was required for inference.

LoRA fine-tuning (Condition 4) was conducted on a JupyterHub instance with
an NVIDIA A100 MIG \texttt{2g.20gb} partition (20\,GB GPU memory,
32\,CPU cores, 64\,GB RAM).
Models were loaded in 4-bit NF4 quantization using QLoRA
\citep{dettmers2023qlora}, reducing peak GPU memory to approximately 8\,GB
during forward passes and 14\,GB during training.

\subsection{Training Details}

\begin{table}[h]
\centering
\caption{LoRA fine-tuning hyperparameters.}
\label{tab:ft_config}
\begin{tabular}{ll}
\toprule
Parameter & Value \\
\midrule
Base model              & mistralai/Mistral-7B-Instruct-v0.2 \\
Quantization            & 4-bit NF4 (QLoRA) \\
LoRA rank ($r$)         & 8 \\
LoRA alpha ($\alpha$)   & 16 \\
Target modules          & \texttt{q\_proj}, \texttt{v\_proj} \\
Training epochs         & 10 \\
Batch size              & 1 (gradient accumulation steps: 4) \\
Learning rate           & $2 \times 10^{-4}$ \\
LR scheduler            & Cosine \\
Warmup steps            & 5 \\
Precision               & bfloat16 \\
Optimizer               & paged\_adamw\_8bit \\
Max sequence length     & 2048 tokens \\
\bottomrule
\end{tabular}
\end{table}

\subsection{Compute Budget}

Table~\ref{tab:compute} summarises total compute usage across all experiments.

\begin{table}[h]
\centering
\caption{Total compute expenditure across all experiments.}
\label{tab:compute}
\begin{tabular}{lrrr}
\toprule
Experiment & Units & Time per unit & Total \\
\midrule
\multicolumn{4}{l}{\textit{Fine-tuning (local GPU)}} \\
LOO folds (17 total)     & 17 folds & $\approx$\,480\,s & $\approx$\,2.3\,h \\
\midrule
\multicolumn{4}{l}{\textit{Inference (Cloudflare Workers AI API)}} \\
Zero-shot (5 models × 17 cases)  & 85 calls  & 10--60\,s & $\approx$\,25\,min \\
Few-shot  (5 models × 17 cases)  & 85 calls  & 15--90\,s & $\approx$\,40\,min \\
RAG       (5 models × 17 cases)  & 85 calls  & 10--60\,s & $\approx$\,25\,min \\
LOO fine-tuned (17 cases)         & 17 calls  & 20--60\,s & $\approx$\,10\,min \\
LLM-as-judge (6 × 17)            & 102 calls & 3--5\,s   & $\approx$\,8\,min \\
\midrule
\multicolumn{4}{l}{\textit{Embedding (BGE-M3)}} \\
Chunk embeddings (2,838 chunks)  & 142 batches & $<$1\,s  & $\approx$\,3\,min \\
Query embeddings (RAG, 17 cases) & 17 calls    & $<$1\,s  & $<$1\,min \\
\midrule
\textbf{Total wall-clock time}   &  &  & $\approx$\,3.5\,h \\
\bottomrule
\end{tabular}
\end{table}

Cloudflare Workers AI was used for all inference and embedding calls during
the open beta period, during which the service was free of charge.
The GPU instance for fine-tuning was provided by institutional research
computing infrastructure.
The dataset construction pipeline (scraping, parsing, extraction) required
an additional $\approx$4--6\,hours of CPU time, primarily for PDF parsing
with \texttt{pdfplumber} and structured extraction via the CF Workers AI API
across all 33 records.

\subsection{Reproducibility}

All experiments are fully reproducible from the released code and dataset.
Key version information:
\begin{itemize}
\item \texttt{transformers} 4.40+,
      \texttt{peft} 0.10+,
      \texttt{bitsandbytes} 0.43+
\item \texttt{pdfplumber} 0.11
\item \texttt{requests} 2.32
\item \texttt{beautifulsoup4} 4.14
\item Cloudflare Workers AI API (\texttt{BGE-M3},
      \texttt{llama-3.3-70b-instruct-fp8-fast},
      \texttt{mistral-7b-instruct-v0.2-lora})
\end{itemize}
Random seeds are fixed (seed 42) for all sampling operations.
Fine-tuning results may vary by $\pm$0.01 F1 across hardware due to
non-deterministic bfloat16 accumulation.

%% ─────────────────────────────────────────────────────────────────────────────

\section{Dataset Card}
\label{app:datasheet}

Following \citet{gebru2021datasheets}, we provide a structured summary of the
dataset.

\textbf{Motivation.}
The dataset was created to enable benchmarking of LLMs on transportation
impact analysis mitigation recommendation generation, a novel application
domain not previously studied in NLP or AI.

\textbf{Composition.}
33 records from California (23) and Massachusetts (10), spanning 2012--2024.
Each record contains a structured project input (project type, location,
description, existing LOS conditions) and ground-truth mitigation measures
extracted from real EIR transportation chapters.
17 of 33 records additionally contain structured issues (failing intersections
with LOS grades), enabling the conditioned generation task.

\textbf{Collection process.}
Documents were collected programmatically from two public portals:
CEQAnet (server-rendered HTML, scraped by date) and Massachusetts MEPA
(REST API, reverse-engineered from browser network traffic).
Ground truth was extracted using Cloudflare Workers AI
(\texttt{@cf/meta/llama-3.3-70b-instruct-fp8-fast}) via structured prompts.

\textbf{Preprocessing.}
PDF text was extracted using \texttt{pdfplumber}.
Section boundaries were detected via regular expression patterns for both
standard TIA headers and California EIR-specific variants.
A leakage check verified that no ground-truth language appeared in input fields.

\textbf{Uses.}
The dataset is intended for benchmarking LLMs on domain-specific structured
generation tasks in transportation planning.
It should not be used to evaluate models for deployment in actual TIA
preparation without additional validation by licensed traffic engineers.

\textbf{Distribution.}
The dataset construction pipeline and benchmark code are available at
\url{https://github.com/[author]/tia-llm-benchmark}.
Source documents are publicly available from CEQAnet and MEPA under
California and Massachusetts public records law respectively.

\textbf{Maintenance.}
The dataset reflects LOS-based analysis methodology (pre-2023 California,
all Massachusetts records).
It does not cover VMT-based analysis required under California SB~743.

%% ─────────────────────────────────────────────────────────────────────────────

\section{Extended Qualitative Error Analysis}
\label{app:qualitative}

This appendix provides extended prediction examples for the zero-shot condition
(Llama-3.3-70B), illustrating the three failure modes identified in
Section~\ref{sec:results}.
For each case the table shows the failing intersection provided as input
(with LOS grade), the ground-truth recommendation from the source document,
the model prediction, and the failure classification.
Bold text highlights the key mismatch.

\subsection*{Failure Mode 1: Correct Location, Wrong Intervention Type}

The model names the right intersection but prescribes a generic engineering
fix rather than the planning-level obligation stated in the ground truth.
This is common for planning-scope documents (EIRs) where the mitigation
is a design standard obligation rather than a specific physical improvement.

\begin{table}[h]
\centering
\caption{Failure Mode 1: correct location, wrong type (zero-shot, Llama-3.3-70B).}
\begin{tabular}{p{3.0cm}p{4.2cm}p{4.2cm}}
\toprule
Input intersection & Ground truth & Prediction \\
\midrule
Lafayette St \& El Camino Real (LOS D, PM peak) &
\texttt{access\_modification}: ``Driveways of future developments shall meet minimum width requirements per City Code'' &
\texttt{signal\_timing\_optimization}: ``Optimize signal timing to extend green time on El Camino Real northbound'' \\
\addlinespace
Holly Oak Drive site access (LOS E, AM peak) &
\texttt{access\_modification}: ``Restrict outbound left turns to right-turn only via a raised median nose'' &
\texttt{new\_traffic\_signal}: ``Install a traffic signal at Holly Oak Drive to manage increased project movements'' \\
\bottomrule
\end{tabular}
\end{table}

\subsection*{Failure Mode 2: Correct Type, Wrong Sub-Location}

The model identifies the right category of improvement at the right arterial
but targets the signalized main-line intersection rather than the correct
sub-location (a site driveway, a ramp terminal, or a secondary intersection).
This sub-location error is sufficient to score as a false positive under
both Tier-1 and Tier-2 matching when the specific location token does not overlap.

\begin{table}[h]
\centering
\caption{Failure Mode 2: correct type, wrong sub-location (zero-shot, Llama-3.3-70B).}
\begin{tabular}{p{3.0cm}p{4.2cm}p{4.2cm}}
\toprule
Input intersection & Ground truth & Prediction \\
\midrule
Harvill Av. \& Perry St. (LOS F, PM peak) &
\texttt{access\_modification}: ``Install stop control and 100-ft WB left-turn lane at \textbf{Driveway 1}'' &
\texttt{add\_turn\_lane}: ``Add dedicated \textbf{northbound} left-turn lane on Harvill Avenue at Perry Street'' \\
\addlinespace
I-215 NB Ramps \& Ramona Exwy. (LOS E, PM peak) &
\texttt{add\_turn\_lane}: ``Add second EBL lane and third EB/WB through lane at \textbf{I-215 SB Ramps}'' &
\texttt{add\_turn\_lane}: ``Add a westbound left-turn lane at \textbf{I-215 NB Ramps and Ramona Expressway}'' \\
\bottomrule
\end{tabular}
\end{table}

\subsection*{Failure Mode 3: Project-Type Mismatch}

For large transit-oriented or institutional developments, ground-truth
mitigations often involve multimodal corridor improvements, street extensions,
or bicycle infrastructure.
The model defaults to intersection-level signal or turn-lane fixes regardless
of project type, reflecting a bias toward the most frequent mitigation pattern
in general TIA documents.

\begin{table}[h]
\centering
\caption{Failure Mode 3: project-type mismatch (zero-shot, Llama-3.3-70B).}
\begin{tabular}{p{3.0cm}p{4.2cm}p{4.2cm}}
\toprule
Input intersection & Ground truth & Prediction \\
\midrule
Potter St./ Third St. (LOS E, 22 failing intersections) &
\texttt{pedestrian\_improvement}: ``Construct sidewalk-level separated bicycle lanes on Potter Street'' &
\texttt{new\_traffic\_signal}: ``Install traffic signal at Potter Street / Third Street'' \\
\addlinespace
Grove St. at I-95 SB Ramps (LOS F, PM peak) &
\texttt{roundabout\_installation}: ``Conversion to a one-lane roundabout at Grove Street / I-95 SB Ramps'' &
\texttt{signal\_timing\_optimization}: ``Retiming of I-95 ramp terminal signal to provide additional green for EB approach'' \\
\bottomrule
\end{tabular}
\end{table}

\subsection*{Comparison Across Conditions}

Table~\ref{tab:qualitative_rag} shows the same three cases under the few-shot
condition (Llama-3.3-70B), illustrating partial improvement in sub-location
specificity when demonstrations are provided.

\begin{table}[h]
\centering
\caption{Few-shot condition on the same cases: partial improvement in
location specificity, project-type mismatch persists.}
\label{tab:qualitative_rag}
\begin{tabular}{p{3.0cm}p{4.2cm}p{4.2cm}}
\toprule
Case / Input & Ground truth & Few-shot prediction \\
\midrule
Harvill Av. site driveway (Mode 2) &
\texttt{access\_modification} @ Driveway 1 &
\texttt{access\_modification} @ ``project site driveway on Harvill Avenue''
(\emph{closer but imprecise}) \\
\addlinespace
Grove St. I-95 SB Ramps (Mode 3) &
\texttt{roundabout\_installation} @ Grove St. I-95 SB Ramps &
\texttt{signal\_timing\_optimization} @ Grove St. I-95 SB Ramps
(\emph{location correct, type wrong}) \\
\addlinespace
Lafayette St. (Mode 1) &
\texttt{access\_modification} @ Future developments &
\texttt{access\_modification} @ ``project driveways''
(\emph{type now correct; location still generic}) \\
\bottomrule
\end{tabular}
\end{table}

The few-shot condition shows that demonstrations help models adopt the
correct intervention vocabulary (Mode 1 now yields \texttt{access\_modification}
rather than \texttt{signal\_timing\_optimization}) but do not resolve the
sub-location precision required for Mode 2 or the project-type vocabulary
gap that produces Mode 3 errors.

%% ─────────────────────────────────────────────────────────────────────────────
\bibliographystyle{elsarticle-num-names}
\bibliography{references}
%% ─────────────────────────────────────────────────────────────────────────────

\end{document}
